\documentclass[11pt]{article}

\usepackage[letterpaper,margin=1in]{geometry}

\usepackage{akteach}
\usepackage{amsmath}
\usepackage{amssymb}
\usepackage{enumerate}
\usepackage{keystroke}

\usepackage{natbib}
\newcommand{\doi}[1]{\href{http://dx.doi.org/#1}{doi: #1}}

\lstset{
  language=Bash,
  basicstyle=\ttfamily,
  commentstyle=\color{blue},
  stringstyle=\color{black},
}

\begin{document}

\akteachheader{High-Performance Scientific Computing (MATH-GA 2011/ CSCI-GA 2945)}%
{Homework Set 3}
\akteachsubheader{Out: September 26, 2012 $\cdot$ Due: October 3, 2012 }

\bigskip
\akteachprobhead{Problem 1: Image + OpenCL warm-up}

Read in two numbers $w$ and $h$ from the command line and allocate three
buffers \texttt{red}, \texttt{green}, \texttt{blue} on the compute
device, each sufficient for $w\times h$ \texttt{unsigned char}s.


Write an OpenCL kernel that fills the buffers above with
linearly blended colors such that
\begin{itemize}
  \item The red channel is zero on the left-hand side of your
    image and 255 (the maximum channel value) on the right-hand side
    of your image.
  \item The blue channel is zero on the top of your
    image and 255 (the maximum channel value) on the bottom
    of your image.
  \item The green channel is zero everywhere.
\end{itemize}
Write your kernel with a work group size of $1\times 1$ and a
global size of $w\times h$.

Use the library routines provided in
\weblink{https://github.com/hpc12/hw3-ppm}{this repository} to
write a \weblink{https://en.wikipedia.org/wiki/Netpbm_format}{PPM}
(``portable pixmap'') image file called \texttt{linear-blend.ppm}.
The file \texttt{test-ppm.c} demonstrates the use of the library.

You may look at the resulting image by typing
\begin{lstlisting}
display linear-blend.ppm
\end{lstlisting}

This command is preinstalled on your virtual machine. If you're not
using the VM, find an image viewer that can display PPM images.

Note that images are conventionally represented in row-major
order. In other words, the $x$ coordinate varies fastest in
memory.

Turn in a main C file \texttt{problem-1/blend.c} along with
a kernel file \texttt{problem-1/linear-blend.cl}.

Also, please make sure to \emph{not} check the generated image files
into git. (You'll kill my server if you do. I'm not kidding.)

\bigskip
\akteachprobhead{Problem 2: Compute a Mandelbrot set}

The \weblink{https://en.wikipedia.org/wiki/Mandelbrot_set}{Mandelbrot
set} is a subset of the complex numbers that has a surprisingly simple
mathematical definition, but a surprisingly complicated mathematical
structure. (Don't know complex numbers?  Don't panic, see below.)

Somewhat formally, it is defined as
\[
  \{c \in \mathbb C: \text{The sequence defined by $z_0=0$,
  $z_n=z_{n-1}^2 +c$ becomes $> 2$ in magnitude for some $n$}\}.
\]
So, in order to determine

If your complex number skills are a bit rusty: A complex number
consists of two real numbers (think \texttt{float}s) $z=(a,b)$.
If we let $c=(x,y)$ and $z_n=(a_n,b_n)$ (for $n\ge 0$), then the
above formula can be expressed as
\begin{align*}
  a_n &= x-b_{n-1}^2+a_{n-1}^2\\
  b_n &= y-2a_{n-1}b_{n-1}
\end{align*}

\begin{enumerate}[a)]
  \item Read in two numbers $w$ and $h$ from the command line and allocate three
    buffers \texttt{red}, \texttt{green}, \texttt{blue} on the compute
    device, each sufficient for $w\times h$ \texttt{unsigned char}s.

  \item Write an OpenCL kernel that carries out the Mandelbrot
    iteration for $c=(x,y)$. Let $x$ vary linearly from
    \texttt{xleft} and \texttt{xright} on the left and right, and
    let $y$ vary linearly from \texttt{ytop} and \texttt{ybottom},
    on the respective ends of your image. Pass these values as
    parameters to your kernel.

    Use the following values for the work you turn in:
    \begin{center}
      \begin{tabular}{|l|l|}
        \hline
        \texttt{xleft} & $-2.13$ \\
        \hline
        \texttt{xright} & 0.77 \\
        \hline
        \texttt{ytop} & 1.3 \\
        \hline
        \texttt{ybottom} & $-1.3$\\
        \hline
      \end{tabular}
    \end{center}
    But feel free to play around.

    For each pixel of the image, run at most \texttt{max\_iter}
    iterations, where that is another parameter to your kernel. Stop
    the iteration when the square of the magnitude of the iterate
    reaches greater than four. For a complex number $z=(a,b)$, the
    square of the magnitude is computed as $|z|^2=a^2+b^2$.

    Store the number of iterations until four is exceeded in all
    channels (red, green, blue) of your image, in such a way that if
    you used \texttt{max\_iter} iterations, the value becomes 255.
    (You may use more `interesting' color maps, too, if you like.)
    Store the resulting image as \texttt{mandelbrot.ppm}.

    Write your kernel with a work group size of $1\times 1$ and a
    global size of $w\times h$.
    \label{part:mbrot-small-wg}

  \item Modify your OpenCL kernel so that it uses a workgroup size of
    $16\times 16$, but make sure that your program can still be run
    with any $w$ and $h$. (I.e. you will have to handle edge cases.)
    Store the resulting image as \texttt{mandelbrot-large-wg.ppm}.
    \label{part:mbrot-large-wg}

  \item Use the \texttt{timing.h} infrastructure from the last two
    homeworks to time the execution of both the small-workgroup and
    the large-workgroup case. Make sure to wait for the compute device
    to finish its job in the right places. Output the performance in
    millions of pixels per second for each case, to three decimal
    places.

  \item Make sure you free/release all your buffers, command queues, host
    memory, and whatever other resources you've used.

    Also make sure that you check for error returns on all functions
    that can fail, including \texttt{malloc} and \texttt{OpenCL}
    interface functions.
\end{enumerate}

Turn in a main C file \texttt{problem-2/mandelbrot.c} along with
kernel files \texttt{problem-2/small-wg.cl} and
\texttt{problem-2/large-wg.cl}, corresponding to parts
\ref{part:mbrot-small-wg}) and \ref{part:mbrot-large-wg}). Make sure that
\texttt{mandelbrot.c} exercises all parts of this assignment as
described above when compiled and run.

Also, please make sure to \emph{not} check the generated image files
into git. (You'll kill my server if you do. I'm not kidding.)

\bigskip
\akteachprobhead{Problem 3: Blur an image}

In this problem, we're going to make an image ``blurry'' by replacing
each pixel's value with a weighted average of its neighboring pixels.
The weighting of the neighbors is given by this picture (which we'll call the
\emph{kernel}) You can find it in the
\weblink{https://github.com/hpc12/hw3-ppm}{PPM repository}
as \texttt{gaussian-kernel.ppm}:
\begin{center}
  \includegraphics[width=6cm]{gaussian-kernel.png}

  \textbf{Figure:} The blurring ``kernel'', an
  image which will provide the weights for our weighted
  average.
\end{center}
This image was made by \texttt{make-gaussian-kernel.c} in the same
repository, which you can feel free to play with.

If you're on the hunt for fancy words, this whole process of `compute
weighted average of neighbors for each pixel' is called a
`\weblink{https://en.wikipedia.org/wiki/Convolution}{convolution}'.
We let
\begin{itemize}
  \item $f(i,j)$ be the pixel value of the unblurry image at position $(i,j)$ with $(0,0)$
    being the upper left,
  \item $W$ and $H$ be the width and height of the unblurry image,
  \item $k$ be the width of the kernel in pixels (which
    we'll assume is even), and
  \item $K(i,j)$ be the value of the red channel of the
    kernel image such that $K(0,0)$ represents the center.
    We'll also let $l=k/2$.
\end{itemize}
With these definitions we'll be computing the `convolved' image $g$ as
\[
  g(i,j)=\frac{1}{\bar K} \sum_{m,n=-l}^{l-1} K(m,n)f(i-m,j-n)
  \qquad
  (l\le i <W-l,l\le j<H-l)
\]
where
\[
  \bar K:=\sum_{m,n=-l}^{l-1} K(m,n).
\]
\begin{enumerate}[a)]
  \item Load \texttt{gaussian-kernel.ppm} into memory. We will only be
    using information from the red channel.

    Precompute $\bar K$ for the kernel you load in (on the host). For
    \texttt{gaussian-kernel.ppm} it should come out to 51548.

  \item Read a file name from the command line and read the image
    given by that file name. You may assume that both width and height
    are divisible by 16. If you do assume this, make sure you check
    for it and print an error if it's not satisfied.

    If you need an image to test with, either use one that you like (a
    family photo, say) and convert it to PPM by writing
    \begin{lstlisting}
    convert my-treasured-family-memory.jpeg test-image.ppm
    \end{lstlisting}
    \texttt{convert} is preinstalled in the virtual machine. You may
    also pass the option \texttt{-geometry WxH} to \texttt{convert}
    if you need to change the size of the image.
  \item Write an OpenCL kernel that directly implements the
    convolution formula above. Make sure to implement the formula
    using \texttt{floats} by converting all values to that type.
    Convert back to \texttt{unsigned char} (the image channel data
    type) on output. Do this for each channel of your image,
    and write the resulting image to \texttt{blurry.ppm}.

    Do not worry about the boundaries. In other words, only consider
    pixels of your image that are at least $l$ pixels away from the
    boundary.

    Use workgroups of size $16\times 16$.
  \item Write an OpenCL kernel that loads the kernel $K$ into local
    memory before it starts work. Make sure to use synchronization as
    needed. Again, carry out the convolution for each channel 
    and write the resulting image to \texttt{blurry-local.ppm}.

    Use workgroups of size $16\times 16$.

  \item Time the execution of both of the above kernels. Make sure you
    wait for the compute device at all the right spots. Print
    performance data in terms of pixels/s for both kernels.
  \item Once again make sure you free/release all your buffers,
    command queues, host memory, and whatever other resources you've
    used.

    Also make sure that you check for error returns on all functions
    that can fail, including \texttt{malloc}, the \texttt{OpenCL}
    interface functions and the image reading/allocation functions.
\end{enumerate}

Turn in a main C file \texttt{problem-3/convolution.c} along with
kernel files \texttt{problem-3/convolution.cl} and
\texttt{problem-3/convolution-local.cl}. Make sure that
\texttt{convolution.c} exercises all parts of this assignment when
compiled and run.

Also, please make sure to \emph{not} check any large image files
into git. (You'll kill my server if you do. I'm not kidding.)

\end{document}
